\documentclass[fontset=none,11pt,a4paper]{ctexart}
\setCJKmainfont{AR PL SungtiL GB}[BoldFont=AR PL UMing CN,AutoFakeBold=2.5]
\setCJKsansfont{AR PL UMing CN}
\setCJKmonofont{AR PL UMing CN}
\usepackage{geometry,booktabs,multirow,hyperref,enumitem,xcolor,longtable}
\geometry{margin=2.2cm}
\title{多模态Agent Post-Training项目\\完整架构文档}
\author{}
\date{\today}

\begin{document}
\maketitle

\begin{abstract}
本文档对多模态Agent Post-Training项目的全部代码文件、数据文件、模型权重进行系统性梳理。对每个文件说明其功能定位、解决的问题、在整体流程中的位置，以及文件之间的协作关系。文档按模块组织---Agent推理核心(agent/)、训练引擎(train/)、评测系统(eval/)、数据构建管线(scripts/)、训练数据(data/)、模型权重(LoRA adapters)、外部图片资源---覆盖从数据准备到模型部署的完整链路。
\end{abstract}

\tableofcontents

\section{项目总览}

\subsection{核心问题}

给定一张图片和一个自然语言问题，训练VLM严格按照工具调用流程完成任务：

\begin{verbatim}
  User(图片+问题) -> vision_parse(OCR/图表/描述)
    -> retrieve_docs(查规则) / python_exec(计算)
    -> final_answer(给出依据) 或 合理拒答
\end{verbatim}

核心约束：每轮对话模型只能输出一个 \texttt{<tool\_call>} 或 \texttt{<final\_answer>}，不得输出裸文本或多个动作。

\subsection{整体架构}

\begin{verbatim}
  数据构建管线(scripts/)        训练引擎(train/)          推理评测(eval/)
  ┌---------------------+    ┌------------------+    ┌------------------+
  | build_agent_train    |    | train_sft_lora   |    | run_vlm_eval     |
  | build_first/second   |--->| train_dpo_lora   |--->| run_seed_eval    |
  | convert_*           |    +------------------+    +------------------+
  | rebuild_seed_*      |             |                      |
  | build_augmented_*   |             v                      v
  | build_refusal_*     |    模型权重(autodl-tmp/)    评测结果(outputs/)
  | build_dpo_*         |    lora_agent_real_v4       vlm_agent_v4_eval.jsonl
  | build_tool_decision |    lora_agent_v7b_clean     vlm_agent_v7b_eval.jsonl
  +---------------------+
           |
           v
    训练数据(data/sft/)
    sft_agent_train_real_v4_enforced_resized.jsonl
    sft_seed.jsonl
    sft_hard_*.jsonl
           |
           v
    Agent推理核心(agent/)
    runtime.py  tools.py  vlm.py  parser.py  prompts.py
\end{verbatim}

\section{Agent推理核心 (agent/)}

\subsection{runtime.py --- 多轮Agent执行引擎}

\textbf{解决的问题}：管理VLM的多轮工具调用循环。接收用户消息，让VLM生成响应，解析tool\_call或final\_answer，执行工具，将工具结果反馈给模型，循环直到final\_answer或达到max\_steps。

\textbf{核心逻辑}：
\begin{enumerate}[nosep]
    \item 加载system prompt + 用户消息（图片+问题）-> messages
    \item 调用VLM生成 -> 解析输出（parser.parse\_model\_output）
    \item 如果是tool\_call -> 执行对应工具 -> 工具结果append到messages -> 回到步骤2
    \item 如果是final\_answer -> 检查enforce\_required\_tools
    \item 如果漏调必要工具 -> 拦截答案，注入纠正提示 -> 回到步骤2
    \item 如果通过 -> 返回answer + trace
\end{enumerate}

\textbf{协作关系}：依赖tools.py（工具执行）、parser.py（输出解析）、prompts.py（系统提示词+纠正提示词模块）、vlm.py（模型推理）。

\textbf{Enforce机制}：通过关键词匹配判断问题类型，如果问题是发票类+未调retrieve\_docs，或图表类+未调python\_exec，拦截final\_answer并注入纠正提示。这是V2从50\% tool hit提升到100\%的关键。

\subsection{tools.py --- 工具实现}

\textbf{解决的问题}：提供vision\_parse（视觉解析）、retrieve\_docs（规则检索）、python\_exec（数学计算）三个工具的真实实现。

\textbf{vision\_parse}：
\begin{itemize}[nosep]
    \item 支持4种模式：ocr、caption、chart\_values、ocr+caption
    \item 真实VLM推理：当\_vlm\_model被注入时，调用Qwen-VL进行视觉推理
    \item Mock回退：当VLM不可用时，根据文件名模式返回预设模拟结果
    \item 画像缓存：对外部数据集使用image+mode级缓存避免重复推理
\end{itemize}

\textbf{retrieve\_docs}：
\begin{itemize}[nosep]
    \item 加载data/docs/下的规则文档（invoice\_rules.txt, chart\_rules.txt, product\_rules.txt）
    \item 通过关键词->文件映射检索对应规则
    \item 支持中英文关键词匹配
\end{itemize}

\textbf{python\_exec}：
\begin{itemize}[nosep]
    \item 正则白名单安全验证：只允许数字、运算符、空格和括号
    \item 在空builtins沙箱中eval执行
\end{itemize}

\textbf{协作关系}：被runtime.py调用。vision\_parse依赖vlm.py的QwenVLModel通过set\_vlm\_model()注入。

\subsection{vlm.py --- VLM模型封装}

\textbf{解决的问题}：封装Qwen2.5-VL模型的加载、推理、LoRA adapter注入、以及mock模式回退。

\textbf{核心类QwenVLModel}：
\begin{itemize}[nosep]
    \item \_\_init\_\_：加载模型和processor，支持adapter注入和mock回退
    \item generate：接收messages列表，返回模型生成的文本
    \item mock\_generate：当VLM不可用时，根据用户文本和图片名推断应输出什么（用于跑通流程）
\end{itemize}

\textbf{协作关系}：被runtime.py（Agent推理）、eval/run\_vlm\_eval.py（评测）使用。当adapter\_name参数非空时，通过PeftModel加载LoRA权重。

\subsection{parser.py --- 模型输出解析}

\textbf{解决的问题}：将VLM的原始文本输出解析为结构化的tool\_call或final\_answer。

\textbf{核心函数parse\_model\_output}：
\begin{enumerate}[nosep]
    \item 正则匹配\texttt{<tool\_call>...</tool\_call>} -> 提取JSON -> 返回\{type: tool\_call, name, args\}
    \item 容错：如果缺闭合标签，尝试提取JSON对象
    \item 正则匹配\texttt{<final\_answer>...</final\_answer>} -> 返回\{type: final\_answer, content\}
    \item 容错：如果缺闭合标签但内容非空，尝试恢复
    \item 以上都不匹配 -> 返回\{type: parse\_error\}
\end{enumerate}

\textbf{协作关系}：被runtime.py和eval系统调用。是格式纪律的最后一道防线。

\subsection{prompts.py --- 提示词模块}

\textbf{解决的问题}：管理系统提示词模板。定义VLM的工具调用格式要求（每轮只能一个tool\_call或final\_answer）。

\section{训练引擎 (train/)}

\subsection{train\_sft\_lora.py --- SFT LoRA训练}

\textbf{解决的问题}：对Qwen2.5-VL基座模型进行LoRA微调，使模型学会多轮工具调用流程。

\textbf{核心机制}：
\begin{itemize}[nosep]
    \item AgentSFTDataset：读取JSONL训练数据，通过expand\_assistant\_turns将每条多轮对话拆分为多个训练样本（每个assistant turn作为一个样本）
    \item QwenVLSFTCollator：对每个样本，计算prefix（目标turn之前的所有内容）和full（包含目标turn），在labels中将prefix位置的token设为-100（不计算loss），只对assistant输出的token计算loss
    \item LoRA配置：r=16/32/64, alpha=2r, target\_modules覆盖所有attention和MLP投影层
\end{itemize}

\textbf{支持的参数}：model\_name、train\_file、output\_dir、num\_epochs、lr、bf16/fp16、gradient\_checkpointing、lora\_r、lora\_alpha、adapter\_name（从已有adapter微调）。

\textbf{协作关系}：输入=scripts/构建的训练数据JSONL，输出=LoRA adapter权重（保存到autodl-tmp/）。

\subsection{train\_dpo\_lora.py --- DPO偏好训练}

\textbf{解决的问题}：在SFT基础上通过偏好数据优化答案质量。但由于格式崩溃问题，三次尝试均失败。

\textbf{核心机制}：
\begin{itemize}[nosep]
    \item DPOTrainer继承Trainer，重写compute\_loss
    \item 加载ref\_model（SFT adapter merge后冻结）和policy\_model（新LoRA）
    \item compute\_log\_prob：只计算最后一条assistant消息的log-prob
    \item DPO loss = -log(sigmoid(beta * (chosen\_lp - ref\_chosen\_lp - rejected\_lp + ref\_rejected\_lp)))
\end{itemize}

\textbf{失败原因}：DPO只优化last-turn，不保护多轮格式纪律。详细分析见前文问题4。

\textbf{协作关系}：输入=DPO偏好数据JSONL + SFT adapter路径，输出=DPO LoRA权重。

\section{评测系统 (eval/)}

\subsection{run\_vlm\_eval.py --- VLM评测主脚本}

\textbf{解决的问题}：加载VLM模型（带LoRA adapter），对50条评测样本逐一运行Agent推理，记录完整的工具调用trace和最终答案，计算多项指标。

\textbf{核心指标}：
\begin{itemize}[nosep]
    \item Tool Hit：gold\_tools是pred\_tools的子集
    \item Tool Order Hit：gold\_tools顺序一致
    \item Format Valid：无parse\_error
    \item Finished：正常结束（未超max\_steps）
    \item Answer Keyword Hit：gold\_keywords全部出现在答案中
    \item Refusal Correct：should\_refuse与is\_refusal一致
\end{itemize}

\textbf{协作关系}：输入=model\_name + adapter\_name + eval\_path，输出=JSONL评测结果。

\subsection{run\_seed\_eval.py --- 早期评测}

\textbf{解决的问题}：项目早期的简化版评测，不依赖真实VLM。用于快速验证Agent逻辑。

\section{数据构建管线 (scripts/)}

\subsection{核心数据构建}

\subsubsection{build\_first\_tool\_sft.py}

\textbf{解决的问题}：从完整的seed数据（sft\_seed.jsonl）中提取"用户消息->第一个tool\_call"的训练样本。只保留到第一个工具调用为止，用于训练模型的第一步决策。

\textbf{协作关系}：输入=sft\_seed.jsonl -> 输出=sft\_first\_tool\_seed.jsonl。

\subsubsection{build\_second\_tool\_sft.py}

\textbf{解决的问题}：从seed数据中提取从开头到第二个工具调用（retrieve\_docs或python\_exec）的训练样本。用于增强模型对多工具链的学习。

\textbf{协作关系}：输入=sft\_seed.jsonl + sft\_hard\_full\_trace.jsonl -> 输出=sft\_second\_tool\_retrieve\_docs.jsonl。

\subsubsection{build\_sft\_train.py}

\textbf{解决的问题}：合并first\_tool和full\_trace数据，去重和验证。产出最终的first\_tool\_train和full\_trace\_train数据。

\textbf{协作关系}：输入=first\_tool\_seed + hard\_first\_tool -> first\_tool\_train；输入=seed + hard\_full\_trace + hard\_format -> full\_trace\_train。

\subsubsection{build\_agent\_train\_v3.py}

\textbf{解决的问题}：合并所有SFT数据源（first\_tool + full\_trace + second\_tool）为一个统一的训练集。

\subsubsection{build\_agent\_train\_real\_v3.py}

\textbf{解决的问题}：合并SFT数据+外部真实数据集（ChartQA、DocVQA、CORD），控制上采样倍率和recovery变体数量，产出最终训练集。

\textbf{核心参数}：
\begin{itemize}[nosep]
    \item external\_sample\_ratio：外部数据采样比例（0.5=50\%）
    \item python\_exec\_repeat/retrieve\_docs\_repeat：上采样倍率
    \item recovery\_repeat：recovery变体数量
\end{itemize}

\subsection{外部数据集转换}

\subsubsection{convert\_chartqa\_to\_sft.py}

\textbf{解决的问题}：将HuggingFace的ChartQA数据集转换为Agent SFT格式。每个图表QA对包装为vision\_parse(chart\_values) -> [python\_exec] -> final\_answer的tool调用链。

\textbf{协作关系}：通过--use\_real\_vlm标志调用real\_vision\_parse.py进行真实VLM推理。

\subsubsection{convert\_docvqa\_to\_sft.py}

\textbf{解决的问题}：将DocVQA数据集转换为Agent SFT格式。文档QA包装为vision\_parse(ocr) -> [retrieve\_docs] -> final\_answer。

\subsubsection{convert\_cord\_to\_sft.py}

\textbf{解决的问题}：将CORD收据数据集转换为Agent SFT格式。每张收据生成4个QA对，包装为vision\_parse(ocr) -> retrieve\_docs -> final\_answer。

\subsection{真实VLM数据重建}

\subsubsection{real\_vision\_parse.py}

\textbf{解决的问题}：提供共享的真实VLM vision\_parse模块。封装Qwen-VL模型的加载和视觉推理，供所有转换脚本使用。

\subsubsection{rebuild\_seed\_with\_real\_vlm.py}

\textbf{解决的问题}：用真实Qwen-VL模型重建任意SFT JSONL文件中的vision\_parse结果。将mock/占位文本替换为真实VLM输出。支持image+mode级缓存。

\textbf{协作关系}：处理seed数据(46张图)+ChartQA(500张)+DocVQA(200张)+CORD(51张)，共797张图片。

\subsection{数据增强与修正}

\subsubsection{build\_augmented\_v5\_data.py}

\textbf{解决的问题}：从seed数据生成多样化问题的增强训练数据。包括：多样问法模板（每个样本×5-7种问法）、推理链前缀、纠错轨迹。

\subsubsection{build\_hard\_cases\_from\_eval.py}

\textbf{解决的问题}：从评测错误中自动生成SFT修正样本和DPO偏好对。错误类型包括：missing\_tool\_call、wrong\_tool\_selection、over\_refusal、should\_refuse\_but\_answered等。

\textbf{协作关系}：输入=tag\_vlm\_errors.py的输出（tagged eval结果），输出=hard\_first\_tool + hard\_format + hard\_full\_trace + dpo\_hard\_cases。

\subsubsection{build\_dpo\_multiturn\_v4.py}

\textbf{解决的问题}：构建多轮DPO偏好对。与单轮DPO不同，prompt包含完整的工具调用上下文，只优化最后一步决策。

\subsubsection{build\_refusal\_samples.py}

\textbf{解决的问题}：用真实VLM（7B）对logo图片做vision\_parse，将"数字/乱码->拒答"的正确流程生成为SFT训练样本。

\subsubsection{build\_tool\_decision\_samples.py}

\textbf{解决的问题}：生成教模型"何时不调工具"和"何时调工具后拒答"的训练样本。包括带推理前缀的直接回答样本和工具调用+拒答样本。

\subsection{诊断与分析}

\subsubsection{tag\_vlm\_errors.py}

\textbf{解决的问题}：对评测结果进行自动错误标注。将每个样本分类为ok、missing\_tool\_call、wrong\_tool\_selection、answer\_keyword\_miss等类别。

\subsubsection{analyze\_sft\_train.py, analyze\_vlm\_eval.py, analyze\_error\_tags.py}

\textbf{解决的问题}：训练数据和评测结果的分析工具。统计工具调用分布、类别分布、错误标签分布。

\subsubsection{add\_eval\_category.py, check\_seed\_data.py}

\textbf{解决的问题}：eval数据集预处理（添加category标签）和seed数据完整性检查。

\section{训练数据 (data/)}

\subsection{SFT训练数据}

\begin{longtable}{lp{9cm}}
\toprule
文件 & 说明 \\
\midrule
sft\_seed.jsonl & 50条种子数据。手动构造的mock vision\_parse结果的完整tool调用链 \\
sft\_first\_tool\_train.jsonl & 84条。用户消息->第一个tool\_call（训练第一步决策） \\
sft\_full\_trace\_train.jsonl & 85条。完整的多轮工具调用轨迹 \\
sft\_hard\_first\_tool.jsonl & 34条。从评测错误中生成的first\_tool修正样本 \\
sft\_hard\_full\_trace.jsonl & 34条。从评测错误中生成的full\_trace修正样本 \\
sft\_hard\_format.jsonl & 1条。格式修正样本 \\
sft\_agent\_train\_real\_v4\_enforced\_resized.jsonl & 1133条。V4/V7B训练数据（核心）。合并seed+hard+外部数据集，resize图片版本 \\
sft\_agent\_train\_v7b\_final.jsonl & 1136条。V7B-final训练数据（V4数据+3条refusal样本） \\
sft\_refusal\_boundary\_v7.jsonl & 9条。用7B真实VLM生成的防幻觉样本 \\
sft\_tool\_decision\_v8.jsonl & 53条。工具决策训练样本（直接回答+refusal+tool） \\
\bottomrule
\end{longtable}

\subsection{评测数据}

\begin{longtable}{lp{9cm}}
\toprule
文件 & 说明 \\
\midrule
eval\_dev.jsonl & 50条。评测集（不含category标签） \\
eval\_dev\_with\_category.jsonl & 50条。带category标签的评测集（使用中） \\
eval\_test.jsonl & 空文件（测试集未构建） \\
\bottomrule
\end{longtable}

\subsection{DPO偏好数据}

\begin{longtable}{lp{9cm}}
\toprule
文件 & 说明 \\
\midrule
dpo\_seed.jsonl & 30条。种子DPO偏好对 \\
dpo\_hard\_cases.jsonl & 28条。从V3评测错误生成的DPO硬case \\
dpo\_train\_v4\_mt.jsonl & 13条。多轮DPO偏好对（V4错误修正） \\
dpo\_visionarena.jsonl & 空文件（未使用） \\
\bottomrule
\end{longtable}

\subsection{规则文档 (data/docs/)}

\begin{longtable}{lp{9cm}}
\toprule
文件 & 说明 \\
\midrule
invoice\_rules.txt & 发票必填字段规则：Invoice No, Date, Amount, Tax ID \\
chart\_rules.txt & 图表计算规则：同比增长率 = (当前值-基期值)/基期值×100\% \\
product\_rules.txt & 商品图文一致性规则：颜色、品类、品牌任意冲突则不一致 \\
\bottomrule
\end{longtable}

\subsection{图片资源 (data/images/)}

\textbf{解决的问题}：提供训练和评测用的视觉输入。

\begin{longtable}{lp{9cm}}
\toprule
类别 & 文件 \\
\midrule
图表 (12张) & chart\_001.png -- chart\_012.png \\
发票 (17张) & invoice\_001.jpg -- invoice\_017.jpg \\
品牌Logo (7张) & logo\_001.jpg -- logo\_007.jpg \\
商品图片 (17张) & product\_001.jpg -- product\_017.jpg \\
\bottomrule
\end{longtable}

共53张图片，覆盖4个任务类型。Seed数据引用其中46张。

\section{模型权重 (autodl-tmp/multimodal-agent-lora/)}

\begin{longtable}{lp{7cm}p{2.5cm}}
\toprule
模型 & LoRA配置 & 状态 \\
\midrule
lora\_agent\_real\_v2 & 3B, r=16, mock数据, 446条 & 基准线 \\
lora\_agent\_real\_v4 & 3B, r=16, 真实VLM数据, 1133条 & \textbf{3B最优} \\
lora\_agent\_v7b\_clean & 7B, r=32, 真实VLM数据, 1133条 & \textbf{7B最优} \\
\bottomrule
\end{longtable}

\section{外部图片资源 (autodl-tmp/multimodal-agent-data/)}

\begin{longtable}{lp{5cm}p{5cm}}
\toprule
目录 & 内容 & 说明 \\
\midrule
chartqa/images\_resized/ & 500张 & ChartQA图表，已resize到1024px \\
docvqa/images\_resized/ & 200张 & DocVQA文档，已resize \\
cord/images\_resized/ & 51张 & CORD收据，已resize（原图203条，51张去重） \\
\bottomrule
\end{longtable}

原始大图已删除（OOM问题，仅保留resize版本）。

\section{完整工作流}

\subsection{训练数据构建流程}

\begin{verbatim}
1. sft_seed.jsonl (50条, 手动构造)
2. build_first_tool_sft.py -> sft_first_tool_seed.jsonl
3. build_second_tool_sft.py -> sft_second_tool_retrieve_docs.jsonl
4. build_sft_train.py -> sft_first_tool_train.jsonl + sft_full_trace_train.jsonl
5. convert_chartqa_to_sft.py --use_real_vlm -> sft_chartqa_real_vlm.jsonl
6. convert_docvqa_to_sft.py --use_real_vlm -> sft_docvqa_real_vlm.jsonl
7. convert_cord_to_sft.py --use_real_vlm -> sft_cord_real_vlm.jsonl
8. rebuild_seed_with_real_vlm.py -> sft_seed_real_vlm.jsonl
9. build_agent_train_real_v3.py -> sft_agent_train_real_v4_enforced.jsonl
10. resize外部图片 -> *_resized/ 目录
11. 更新JSONL路径 -> sft_agent_train_real_v4_enforced_resized.jsonl
\end{verbatim}

\subsection{训练流程}

\begin{verbatim}
train_sft_lora.py
  --model_name /path/to/Qwen-VL
  --train_file sft_agent_train_real_v4_enforced_resized.jsonl
  --output_dir /path/to/lora_adapter
  --bf16 --lora_r 16 --gradient_accumulation_steps 8
  -> LoRA adapter权重
\end{verbatim}

\subsection{评测流程}

\begin{verbatim}
run_vlm_eval.py
  --model_name /path/to/Qwen-VL
  --adapter_name /path/to/lora_adapter
  --enforce_required_tools True/False
  -> JSONL评测结果 (trace + metrics)
\end{verbatim}

\subsection{推理流程}

\begin{verbatim}
QwenVLModel(model_name, adapter_name).generate(messages)
  -> runtime.run(image, question)
    -> 多轮循环: VLM生成 -> parser解析 -> 工具执行 -> 结果反馈
    -> {answer, trace, status}
\end{verbatim}

\section{文件协作关系图}

\begin{verbatim}
                    ┌----------------------------------+
                    |         data/images/ (53张)       |
                    |    chart_*.png  invoice_*.jpg    |
                    |    logo_*.jpg   product_*.jpg    |
                    +----------+-----------------------+
                               | 视觉输入
                    ┌----------v-----------------------+
                    |      agent/tools.py              |
                    |  vision_parse(retrieve_docs)     |
                    |  python_exec                     |
                    +----------+-----------------------+
                               |
    ┌--------------------------┼--------------------------+
    |                          |                          |
    v                          v                          v
agent/vlm.py           agent/runtime.py           agent/parser.py
(模型加载)              (多轮循环)                 (输出解析)
    |                      |                          |
    +----------------------┼--------------------------+
                           |
              ┌------------v------------+
              |    train/               |
              |  train_sft_lora.py     |---> LoRA权重
              |  train_dpo_lora.py     |
              +------------+------------+
                           |
              ┌------------v------------+
              |    eval/                |
              |  run_vlm_eval.py       |---> 评测结果
              |  run_seed_eval.py      |
              +------------+------------+
                           |
              ┌------------v------------+
              |    scripts/             |
              |  数据构建 数据增强       |---> 训练数据JSONL
              |  错误分析 评测标注       |
              |  DPO构建 真实VLM重建    |
              +-------------------------+
\end{verbatim}

\section{总结}

本项目共包含\textbf{约90个文件}，按功能分为6个模块：

\begin{enumerate}[nosep]
    \item \textbf{Agent推理核心}(6个py)：runtime.py、tools.py、vlm.py、parser.py、prompts.py、rule\_agent.py。负责多轮Agent执行、工具调用、输出解析。
    \item \textbf{训练引擎}(2个py)：train\_sft\_lora.py、train\_dpo\_lora.py。负责SFT和DPO的LoRA训练。
    \item \textbf{评测系统}(2个py)：run\_vlm\_eval.py、run\_seed\_eval.py。负责模型评测和指标计算。
    \item \textbf{数据构建管线}(25个py)：覆盖外部数据集转换、真实VLM数据重建、数据增强、错误分析、偏好对构建、工具决策样本生成等全流程。
    \item \textbf{训练数据}(~35个jsonl/txt)：SFT数据、DPO偏好对、评测数据、规则文档。
    \item \textbf{模型与图片}(~850个文件)：3个LoRA adapter、2个基座模型（3B+7B）、53张本地图片+751张外部resize图片。
\end{enumerate}

所有模块围绕一个核心目标协作：\textbf{从数据构建->SFT训练->评测反馈->数据修正，形成闭环迭代}。关键瓶颈不在于某个模块，而在于SFT优化范式的根本局限---这一发现已记录于前文分析中，后续RL/GRPO将是突破路径。

\end{document}
